\chapter{Hơn và Kém}
\markboth{Hơn và Kém}{More and Less}

\section{Một người luôn muốn hơn người nên chẳng khi nào thấy đủ.}
\begin{truyen}{Một người luôn muốn hơn người nên chẳng khi nào thấy đủ.}{Comparison can bankrupt peace.}
\chuhoa{C}{ứ mỗi lần đạt được một điều, anh lại nhìn quanh xem ai đang hơn mình nữa. Vì thế niềm vui chưa kịp ở lại đã bị so sánh đẩy đi. Sống bằng tâm thế phải hơn người khác khiến con người rất khó ở yên với chính mình.}
\textit{(Each time he achieved something, he immediately looked around to see who still had more. Joy never had time to stay because comparison kept pushing it away. Living with the need to surpass others makes it very hard to rest within yourself.)}
\end{truyen}
\begin{baihoc}
Khi cái hơn trở thành thước đo chính, bình an thường là thứ bị kém đi đầu tiên.
\textit{(When being ahead becomes the main measure, peace is often the first thing to diminish.)}
\end{baihoc}
\begin{ghinhoanh}
\textbf{Short English to remember:}

Comparison steals enough.
\end{ghinhoanh}
\begin{tuhocnhanh}
1. Em đang so sánh mình với ai nhiều nhất? \textit{(Whom are you comparing yourself with most often?)}

2. Nếu ngưng cuộc đua đó, em sẽ nhẹ hơn điều gì? \textit{(If you stopped that race, what would become lighter?)}
\end{tuhocnhanh}
\ngancach

\section{Một học sinh kém môn này nhưng lại có năng lực khác bị che khuất.}
\begin{truyen}{Một học sinh kém môn này nhưng lại có năng lực khác bị che khuất.}{Less in one area does not mean less in all.}
\chuhoa{E}{m học Toán chậm nên thường bị xem là học lực yếu. Nhưng khi lớp làm dự án, em lại viết kịch bản tốt, nói chuyện duyên, và kết nối nhóm rất ổn. Có người bị gọi là kém chỉ vì người khác đang nhìn họ qua một khe cửa quá hẹp.}
\textit{(The student learned math slowly and was often labeled weak. Yet when the class did a project, this same student wrote strong scripts, spoke gracefully, and connected the team well. Some people are called lacking only because others are looking through a doorway that is too narrow.)}
\end{truyen}
\begin{baihoc}
Kém ở một thước đo chưa nói hết giá trị của một con người.
\textit{(Being less on one measure does not define a whole person.)}
\end{baihoc}
\begin{ghinhoanh}
\textbf{Short English to remember:}

One weakness is not a whole identity.
\end{ghinhoanh}
\begin{tuhocnhanh}
1. Em có đang đóng khung bản thân bằng một điểm yếu không? \textit{(Are you defining yourself by one weakness?)}

2. Năng lực nào của em chưa được nhìn đúng mức? \textit{(Which of your strengths has not been seen clearly enough?)}
\end{tuhocnhanh}
\ngancach

\section{Hơn người về tiền nhưng kém người về cách sống.}
\begin{truyen}{Hơn người về tiền nhưng kém người về cách sống.}{More possession does not mean more worth.}
\chuhoa{A}{nh thành đạt sớm, tiêu tiền thoải mái và luôn muốn cho người khác thấy mình ở mức cao hơn. Nhưng cách anh cư xử với cha mẹ, với nhân viên, với những người yếu thế lại làm người ta khó kính trọng. Có cái hơn rất sáng ở bề mặt nhưng lại lộ ra một cái kém lớn hơn bên trong.}
\textit{(He succeeded early, spent freely, and always wanted others to see that he lived above them. Yet the way he treated parents, employees, and the vulnerable made respect difficult. Some forms of having more shine on the surface while exposing a deeper inner lack.)}
\end{truyen}
\begin{baihoc}
Hơn về sở hữu không tự động thành hơn về nhân cách.
\textit{(Being ahead in possession does not automatically mean being ahead in character.)}
\end{baihoc}
\begin{ghinhoanh}
\textbf{Short English to remember:}

More wealth is not more character.
\end{ghinhoanh}
\begin{tuhocnhanh}
1. Em thường ngưỡng mộ người khác vì điều gì? \textit{(What do you usually admire in others?)}

2. Em có đang đặt sai thước đo cho cái gọi là hơn không? \textit{(Are you using the wrong measure for what it means to be ahead?)}
\end{tuhocnhanh}
\ngancach

\section{Chấp nhận mình kém ở hiện tại để học thật.}
\begin{truyen}{Chấp nhận mình kém ở hiện tại để học thật.}{Honest humility accelerates growth.}
\chuhoa{K}{hi cô giáo trả bài, em không còn đổ lỗi cho đề khó hay do xui. Em nhìn đúng rằng mình còn kém ở phần đọc hiểu và bắt đầu học lại từng bước. Nhiều người chậm tiến không phải vì kém, mà vì không chịu nhận ra mình đang kém ở đâu.}
\textit{(When the teacher returned the test, the student no longer blamed the hard exam or bad luck. Instead, there was honest recognition of weakness in reading comprehension, followed by step-by-step rebuilding. Many people do not stay behind because they are weak, but because they refuse to name their weakness clearly.)}
\end{truyen}
\begin{baihoc}
Thừa nhận cái kém là khởi đầu của một cái hơn mới.
\textit{(Admitting a weakness is the beginning of a new strength.)}
\end{baihoc}
\begin{ghinhoanh}
\textbf{Short English to remember:}

Name the weakness, start the growth.
\end{ghinhoanh}
\begin{tuhocnhanh}
1. Em đang né nhìn phần nào còn yếu của mình? \textit{(Which weak area of yours are you avoiding?)}

2. Nếu gọi đúng tên cái kém ấy, em sẽ học lại từ đâu? \textit{(If you named that weakness clearly, where would you begin relearning?)}
\end{tuhocnhanh}
\ngancach

\section{Biết mình đủ thay vì cứ đòi mình phải hơn mãi.}
\begin{truyen}{Biết mình đủ thay vì cứ đòi mình phải hơn mãi.}{Enough is sometimes wiser than more.}
\chuhoa{C}{ó người không bỏ phát triển, nhưng thôi không sống bằng cảm giác thiếu triền miên. Họ vẫn học, vẫn làm, vẫn tiến, nhưng không còn xem mình là kẻ thua cuộc chỉ vì chưa hơn ai đó. Khi biết chữ đủ, lòng người bớt rách vì ganh đua.}
\textit{(Some people do not stop growing, but they stop living from a permanent sense of lack. They still learn, work, and move forward, yet no longer see themselves as failures simply because they have not surpassed someone else. When the word enough enters the heart, competition tears it less.)}
\end{truyen}
\begin{baihoc}
Biết đủ không giết ý chí; nó cứu sự bình an.
\textit{(Knowing enough does not kill ambition; it protects peace.)}
\end{baihoc}
\begin{ghinhoanh}
\textbf{Short English to remember:}

Enough protects peace.
\end{ghinhoanh}
\begin{tuhocnhanh}
1. Với em, đủ là gì trong học tập hay cuộc sống? \textit{(For you, what does enough mean in study or life?)}

2. Em có thể tiếp tục cố gắng mà bớt so kè bằng cách nào? \textit{(How can you keep striving while comparing less?)}
\end{tuhocnhanh}
\ngancach
